%-------------------------
% Resume in Latex
% Author : Jake Gutierrez
% Based off of: https://github.com/sb2nov/resume
% License : MIT
%------------------------

\documentclass[letterpaper,11pt]{article}

\usepackage{latexsym}
\usepackage[empty]{fullpage}
\usepackage{titlesec}
\usepackage{marvosym}
\usepackage[usenames,dvipsnames]{color}
\usepackage{verbatim}
\usepackage{enumitem}
\usepackage[hidelinks]{hyperref}
\usepackage{fancyhdr}
\usepackage[french]{babel}
\frenchsetup{StandardLayout=true} % keep the template's bullets and list spacing
\usepackage{tabularx}
\usepackage{iftex}
\ifPDFTeX
  \input{glyphtounicode}
\fi


%----------FONT OPTIONS----------
% sans-serif
% \usepackage[sfdefault]{FiraSans}
% \usepackage[sfdefault]{roboto}
% \usepackage[sfdefault]{noto-sans}
% \usepackage[default]{sourcesanspro}

% serif
% \usepackage{CormorantGaramond}
% \usepackage{charter}


\pagestyle{fancy}
\fancyhf{} % clear all header and footer fields
\fancyfoot{}
\renewcommand{\headrulewidth}{0pt}
\renewcommand{\footrulewidth}{0pt}

% Adjust margins
\addtolength{\oddsidemargin}{-0.5in}
\addtolength{\evensidemargin}{-0.5in}
\addtolength{\textwidth}{1in}
\addtolength{\topmargin}{-.5in}
\addtolength{\textheight}{1.0in}
% Remove the space reserved for the (empty) page header
\setlength{\headheight}{0pt}
\setlength{\headsep}{0pt}

\urlstyle{same}

\raggedbottom
\raggedright
\setlength{\tabcolsep}{0in}

% Sections formatting
\titleformat{\section}{
  \vspace{-4pt}\scshape\raggedright\large
}{}{0em}{}[\color{black}\titlerule \vspace{-5pt}]

% Ensure that generate pdf is machine readable/ATS parsable
% (pdfLaTeX only; XeLaTeX/LuaLaTeX/Tectonic already emit Unicode text)
\ifPDFTeX
  \pdfgentounicode=1
\fi

%-------------------------
% Custom commands
\newcommand{\resumeItem}[1]{
  \item\small{
    {#1 \vspace{-2pt}}
  }
}

\newcommand{\resumeSubheading}[4]{
  \vspace{-2pt}\item
    \begin{tabular*}{\textwidth}[t]{l@{\extracolsep{\fill}}r}
      \textbf{#1} & #2 \\
      \textit{\small#3} & \textit{\small #4} \\
    \end{tabular*}\vspace{-7pt}
}

\newcommand{\resumeSubSubheading}[2]{
    \item
    \begin{tabular*}{\textwidth}{l@{\extracolsep{\fill}}r}
      \textit{\small#1} & \textit{\small #2} \\
    \end{tabular*}\vspace{-7pt}
}

\newcommand{\resumeProjectHeading}[2]{
    \item
    \begin{tabular*}{\textwidth}{l@{\extracolsep{\fill}}r}
      \small#1 & #2 \\
    \end{tabular*}\vspace{-7pt}
}

% One-line heading: title · organisation, location ........ dates
\newcommand{\resumeOneLineHeading}[3]{
  \vspace{-2pt}\item
    \begin{tabular*}{\textwidth}[t]{l@{\extracolsep{\fill}}r}
      \textbf{#1} $\cdot$ #2 & #3 \\
    \end{tabular*}\vspace{-7pt}
}

% Project: name ........ date
%          context | tech stack
\newcommand{\resumeProjectEntry}[4]{
  \vspace{-2pt}\item
    \begin{tabular*}{\textwidth}[t]{l@{\extracolsep{\fill}}r}
      \textbf{#1} & #3 \\
      {\small\textit{#2} $|$ \textit{#4}} & \\
    \end{tabular*}\vspace{-7pt}
}

% Education: school on one line, then one line per degree with dates on the right
\newcommand{\resumeSchoolHeading}[2]{
  \vspace{-2pt}\item
    \begin{tabular*}{\textwidth}[t]{l@{\extracolsep{\fill}}r}
      \textbf{#1} & #2 \\
    \end{tabular*}\vspace{-7pt}
}
% Degree line without a bullet (use inside \resumeSubHeadingListStart)
\newcommand{\resumeDegreeLine}[2]{
  \item\small{#1\hfill#2\vspace{-2pt}}
}
% Bulleted line with dates on the right (Certificates)
\newcommand{\resumeDegreeItem}[2]{
  \item\small{#1\hfill#2\vspace{-2pt}}
}

\newcommand{\resumeSubItem}[1]{\resumeItem{#1}\vspace{-4pt}}

\renewcommand\labelitemii{$\vcenter{\hbox{\tiny$\bullet$}}$}

\newcommand{\resumeSubHeadingListStart}{\begin{itemize}[leftmargin=0pt, label={}]}
\newcommand{\resumeSubHeadingListEnd}{\end{itemize}}
\newcommand{\resumeItemListStart}{\begin{itemize}[labelindent=0.1in, leftmargin=*]}
\newcommand{\resumeItemListEnd}{\end{itemize}\vspace{-5pt}}

%-------------------------------------------
%%%%%%  RESUME STARTS HERE  %%%%%%%%%%%%%%%%%%%%%%%%%%%%


\begin{document}

%----------EN-TETE----------
{\centering
    \textbf{\Huge \scshape Longye Chen} \\ \vspace{2pt}
    \small +33 6 02 53 61 19 $|$ \href{mailto:lchen031205@gmail.com}{\underline{lchen031205@gmail.com}} $|$
    \href{https://www.linkedin.com/in/longye-chen-11ba56309/}{\underline{LinkedIn}} $|$
    \href{https://github.com/CHENLongyeLucien}{\underline{GitHub}}
\par}
\vspace{-10pt}


%-----------FORMATION-----------
\section{Formation}
  \resumeSubHeadingListStart
    \resumeSchoolHeading{Sorbonne Université}{Paris, France}
    \resumeDegreeLine{Licence d'Informatique}{2025 -- aujourd'hui}
      \begin{itemize}[labelindent=0.1in, leftmargin=*, topsep=0pt, partopsep=0pt]
        \resumeItem{\textit{Matières principales :} Science des données, Algorithmique, Structures de données, Bases de données,\\ Développement web, Systèmes, Réseaux}
      \end{itemize}
    \resumeDegreeLine{Licence de Mathématiques et Informatique}{2023 -- 2025}
  \resumeSubHeadingListEnd


%-----------EXPERIENCE-----------
\section{Expérience professionnelle}
  \resumeSubHeadingListStart

    \resumeOneLineHeading
      {Stagiaire Administrateur Systèmes \& Développeur IA}{Lessonia, Bretagne, France}{juin -- sept. 2026}
      \resumeItemListStart
        \resumeItem{Automatisation de l'arrêt ordonné sur coupure de courant avec Eaton IPM (8~hôtes Proxmox, 114~VM, 2~SAN DataCore) : comparaison de 3~solutions, règles par dépendances et arrêt complet en moins de 15~min avec alertes.}
        \resumeItem{Intégration d'un agent IA auto-hébergé (Ollama, Open WebUI) à Proxmox VE et GLPI via des webhooks, générant automatiquement un résumé quotidien de 5 à 10~tickets.}
        \resumeItem{Extension d'un outil Python de requêtes en langage naturel aux bases PostgreSQL et MSSQL, fonctionnant entièrement en local et en lecture seule (modèles Qwen, Gemma).}
      \resumeItemListEnd

    \resumeOneLineHeading
      {Stagiaire Commerce International}{Lessonia, Bretagne, France}{juin -- juil. 2026}
      \resumeItemListStart
        \resumeItem{Gestion de 20+ comptes clients et prospects dans un CRM et suivi des clients internationaux en anglais et chinois.}
        \resumeItem{Prospection de clients potentiels et soutien à l'équipe marketing sur des projets commerciaux.}
      \resumeItemListEnd

  \resumeSubHeadingListEnd


%-----------PROJETS-----------
\section{Projets}
    \resumeSubHeadingListStart
      \resumeProjectEntry
          {\href{https://github.com/CHENLongyeLucien/Feikuai}{\underline{Feikuai}} : Pipeline multi-agents de sécurité des paiements par IA}
          {Hackathon ESSEC $\times$ Kryptosphere}
          {oct. 2026}
          {TypeScript, Solidity, React, Avalanche, Docker}
          \resumeItemListStart
            \resumeItem{Conception en équipe de 3 d'un pipeline multi-agents conforme à DORA, bloquant en amont 97\,\% de 1\,821~attaques, avec 0~paiement frauduleux sur 1\,896~tentatives rejouées on-chain (Avalanche).}
            \resumeItem{Sécurisation de la validation des paiements (143~tests) : développement de l'API, des outils MCP des agents, de la signature EIP-712 et de la gestion des clés côté serveur, et contribution au front-end React.}
          \resumeItemListEnd
      \resumeProjectEntry
          {\href{https://github.com/CHENLongyeLucien/Lumen}{\underline{Lumen}} : Extension de lecteur vidéo personnalisé}
          {Extension Chrome open source}
          {nov. 2026}
          {TypeScript, Manifest V3, Vitest}
          \resumeItemListStart
            \resumeItem{Développement d'une extension Chrome remplaçant le lecteur vidéo natif par un lecteur personnalisé (raccourcis, sous-titres, boucle A-B, dessin, zoom), écrit de zéro, sans dépendance d'exécution, dans un bundle de 17~Ko (gzip).}
            \resumeItem{Couverture de la logique principale par 144~tests automatisés (99\,\% de couverture) ; fonctionne entièrement en local, sans pistage et avec des permissions minimales.}
          \resumeItemListEnd
      \resumeProjectEntry
          {\href{https://github.com/CHENLongyeLucien/Odyssey}{\underline{Odyssey}} : Jeu 3D sur navigateur avec compagnons IA}
          {Hackathon {\NoAutoSpacing\{Tech: Europe\}}}
          {sept. 2026}
          {Gemini, Gradium, Lyria, Vercel}
          \resumeItemListStart
            \resumeItem{Réalisation en équipe de 3 d'un jeu 3D en ligne jouable (30~parcelles constructibles, 4~thèmes de ville, 3~armes, boss final) avec sauvegardes persistantes et prise en charge du clavier et de la manette Xbox.}
            \resumeItem{Intégration de Gemini, Gradium et Lyria pour contrôler les compagnons à la voix ou par texte, en traduisant des ordres en langage naturel en actions validées par les règles du jeu.}
          \resumeItemListEnd
    \resumeSubHeadingListEnd


%-----------COMPETENCES ET LANGUES-----------
\section{Compétences et langues}
 \begin{itemize}[leftmargin=0pt, label={}]
    \small{\item{
     \textbf{Programmation}{ : Python, TypeScript, JavaScript, Java, C, SQL, Solidity, Bash, assembleur RISC-V} \\
     \textbf{Frameworks et bibliothèques}{ : React, Next.js, Node.js, Three.js, FastAPI, Flask} \\
     \textbf{Bases de données}{ : PostgreSQL, MSSQL} \\
     \textbf{Outils et infrastructure}{ : Git, GitHub, Docker, Vitest, Proxmox VE, GLPI, Ollama} \\
     \textbf{Concepts}{ : API REST, agents IA (MCP), extensions Chrome (Manifest V3), tests unitaires} \\
     \textbf{Langues}{ : français (langue maternelle), anglais (C1), chinois (courant)}
    }}
 \end{itemize}


%-----------CERTIFICATIONS-----------
\section{Certifications}
  \resumeItemListStart
    \resumeDegreeItem{IBM SkillsBuild : \href{https://www.credly.com/badges/f42f2621-21d8-478a-9dae-b04f33677714/public_url}{\underline{Artificial Intelligence Fundamentals}}}{sept. 2026}
    \resumeDegreeItem{IBM SkillsBuild : \href{https://www.credly.com/badges/942aa2ea-c60e-4343-9da9-fcaad618f717/public_url}{\underline{Web Development Fundamentals}}}{sept. 2026}
  \resumeItemListEnd


%-------------------------------------------
\end{document}
